\documentclass[11pt]{article}

\usepackage[letterpaper,margin=1in]{geometry}
\usepackage{amsmath,amssymb}
\usepackage{booktabs}
\usepackage{longtable}
\usepackage{array}
\usepackage{xcolor}
\usepackage{listings}
\usepackage{url}
\usepackage{hyperref}
\usepackage{titlesec}
\usepackage{parskip}
\usepackage{fancyhdr}
\usepackage{enumitem}
\usepackage{tikz}
\usetikzlibrary{positioning,arrows.meta,calc,shapes.geometric}

\hypersetup{
  colorlinks=true,
  linkcolor=black,
  urlcolor=blue!50!black,
  citecolor=black,
  pdftitle={LonghornSilicon Software Overview},
  pdfauthor={LonghornSilicon},
  pdfsubject={Compiler co-design entry point}
}

\titleformat{\section}{\normalfont\Large\bfseries}{\thesection.}{0.6em}{}
\titleformat{\subsection}{\normalfont\large\bfseries}{\thesubsection}{0.6em}{}

\pagestyle{fancy}
\fancyhf{}
\lhead{\small LonghornSilicon SW Overview}
\rhead{\small Compiler Co-Design Entry Point}
\cfoot{\thepage}
\renewcommand{\headrulewidth}{0.4pt}

\definecolor{codebg}{gray}{0.97}
\definecolor{codekey}{rgb}{0.13,0.29,0.53}
\definecolor{codecom}{rgb}{0.25,0.5,0.25}
\lstset{
  basicstyle=\ttfamily\small,
  backgroundcolor=\color{codebg},
  frame=single,
  framerule=0.4pt,
  rulecolor=\color{black!30},
  xleftmargin=0.5em,
  xrightmargin=0.5em,
  showstringspaces=false,
  breaklines=true,
  keywordstyle=\color{codekey}\bfseries,
  commentstyle=\color{codecom}\itshape,
  numbers=none
}

\setlist{itemsep=2pt,topsep=3pt}

\begin{document}

%==============================================================================
\thispagestyle{empty}
\begin{center}
\vspace*{1.2in}
{\Huge\bfseries LonghornSilicon Software}\\[0.4em]
{\LARGE\bfseries Compiler Co-Design Overview}\\[1.5em]

\rule{0.7\textwidth}{0.4pt}\\[1em]

\begin{tabular}{rl}
\textbf{Date}:           & 2026-05-14 \\
\textbf{Status}:         & Blocks 1 \& 2 of 4 in software reference \\
\textbf{Audience}:       & Compiler engineers integrating against the chip \\
\textbf{Repository}:     & \texttt{LonghornSilicon/attention-compute-unit} \\
\textbf{Branch}:         & \texttt{compiler-integration-prep} \\
\end{tabular}\\[1em]
\rule{0.7\textwidth}{0.4pt}

\vfill
\begin{minipage}{0.85\textwidth}
\small
\textbf{Purpose.} This document is the orientation map for a compiler
team integrating against the LonghornSilicon chip. It points to the
bit-accurate reference models, the ISA spec, the verification status
of each block, and the open questions we want answered to converge
the v0.2 reference. Read this first; then jump to the specific
references it links.
\end{minipage}

\vspace{0.3in}
\end{center}

\clearpage

\tableofcontents
\clearpage

%==============================================================================
\section{What LonghornSilicon Builds}
\label{sec:what}

A four-block LLM inference accelerator targeting TSMC 16nm FinFET (N16FFC) via the
TSMC University Program (tape-out target Summer 2027):

\begin{enumerate}
\item \textbf{ACU (Attention Compute Unit)} --- decides INT8 vs FP16
      per attention tile and runs the matmul. Composed of:
      \begin{itemize}
      \item Precision Controller --- ratio gate on raw scores
      \item MAC Array --- INT8 and FP16 dispatched matmul
      \end{itemize}
\item \textbf{KV Cache Engine} --- on-die SRAM (spilling to off-chip LPDDR5X) with
      ChannelQuant compression on writes, decompression on reads.
\item \textbf{Token Importance Unit} --- per-token attention-weight
      accumulator driving mixed-precision retention.
\item \textbf{Memory Hierarchy Controller} --- routes 0.8 MB on-die SRAM /
      off-chip LPDDR5X, direct (no separate eDRAM tier).
\end{enumerate}

For compiler co-design, block 1 (the ACU and its two sub-components,
the precision controller and the MAC array) is the first to be fully
specified in software. Blocks 2--4 (KV Cache Engine, Token Importance
Unit, Memory Hierarchy Controller) follow the same template as they
land.

%==============================================================================
\section{Block Coverage in the Software Reference}
\label{sec:coverage}

\begin{table}[ht]
\centering
\footnotesize
\caption{Reference-model status as of \texttt{compiler-integration-prep}.}
\label{tab:coverage}
\setlength{\tabcolsep}{4pt}
\begin{tabular}{@{}p{3.0cm} p{4.0cm} p{1.8cm} p{3.3cm} p{2.3cm}@{}}
\toprule
Block / sub-block           & C++ class                          & C API                 & Python                       & Status \\
\midrule
ACU --- Precision Ctrl.     & \texttt{lhsi::PrecisionController} & \texttt{lhsi\_pc\_*}  & \texttt{PrecisionController} & 143/143 vs RTL \\
ACU --- MAC Array           & \texttt{lhsi::mac::MacArray}       & \texttt{lhsi\_mac\_*} & \texttt{MacArray}            & numpy agreed \\
KV Cache Engine             & ---                                & ---                   & ---                          & in design \\
Token Importance Unit       & ---                                & ---                   & ---                          & not started \\
Memory Hierarchy Controller & ---                                & ---                   & ---                          & not started \\
\bottomrule
\end{tabular}
\end{table}

Each new block follows the same three-language template; see
\path{docs/new_block_blueprint.md} for the per-block scaffolding
pattern.

%==============================================================================
\section{Where Things Live}
\label{sec:layout}

\begin{lstlisting}[basicstyle=\ttfamily\small]
sw/
+-- README.md                         (markdown form of this doc)
+-- reference_model/
    +-- README.md                     (API reference, markdown)
    +-- MAC_ARRAY_DESIGN.md
    +-- Makefile                      (build / test / static / shared / clean)
    +-- precision_controller_ref.{hpp,cpp,py}
    +-- test_precision_controller_ref.{cpp,py}
    +-- mac_array_ref.{hpp,cpp,py}
    +-- test_mac_array_ref.{cpp,py}
    +-- integration_example.py        (end-to-end attention tile)
    +-- example_compiler_use.py       (three sophistication levels)

docs/
+-- isa/precision_controller_isa.pdf  (ISA spec)
+-- mac_array_design.pdf              (this doc family)
+-- reference_model_api.pdf           (formal API reference)
+-- sw_overview.pdf                   (this document)
+-- new_block_blueprint.md            (template for future blocks)
\end{lstlisting}

%==============================================================================
\section{Build and Test (One Line)}
\label{sec:build}

\begin{lstlisting}[basicstyle=\ttfamily\small]
cd sw/reference_model && make test-all
\end{lstlisting}

\noindent
Builds all C++ tests, runs them, runs all Python tests, prints a
final parity statement. Requires Python 3.10+, NumPy, and a C++17
compiler.

Granular targets:

\begin{table}[ht]
\centering
\small
\begin{tabular}{@{}l l@{}}
\toprule
Target          & What it does \\
\midrule
\texttt{test-pc}     & Precision controller C++ test only \\
\texttt{test-mac}    & MAC array C++ test only \\
\texttt{test-py}     & All Python tests \\
\texttt{integration} & End-to-end worked example (printed report) \\
\texttt{shared}      & \texttt{libprecision\_controller\_ref.so} + \texttt{libmac\_array\_ref.so} \\
\texttt{static}      & \texttt{.a} variants \\
\texttt{clean}       & Remove build artifacts \\
\bottomrule
\end{tabular}
\end{table}

%==============================================================================
\section{Recommended Reading Order for a Compiler Engineer}
\label{sec:reading}

\begin{enumerate}
\item \textbf{This document} (top to bottom; you are almost done).
\item \texttt{docs/isa/precision\_controller\_isa.pdf}, especially
      \S4.1 (worked lowering example) and \S4.2 (compiler binding
      patterns: MLIR / TVM / ONNX / custom IR).
\item \texttt{docs/mac\_array\_design.pdf} for the MAC-array design
      decisions (frozen vs.\ open).
\item \texttt{docs/reference\_model\_api.pdf} for the full API
      reference (C++ class, C ABI, Python).
\item \texttt{sw/reference\_model/integration\_example.py} ---
      run it, then read it side-by-side. The shape of that script
      is the shape your codegen should produce.
\end{enumerate}

%==============================================================================
\section{Integration Phases}
\label{sec:phases}

The interface (ISA + reference models) is stable across all four
phases. Only the runtime implementation of the operations changes:
Python function call $\rightarrow$ AXI register write
$\rightarrow$ PCIe MMIO.

\begin{table}[ht]
\centering
\small
\renewcommand{\arraystretch}{1.2}
\caption{Four-phase compiler integration plan.}
\label{tab:phases}
\begin{tabular}{@{}c p{3.0cm} p{4.0cm} p{5.0cm}@{}}
\toprule
Phase & Timeline & Compiler targets & We provide \\
\midrule
0 & now                      & Python + C++ reference models                & Models, ISA, tests on this branch \\
1 & when ZCU102/104 arrives  & AXI-Lite + AXI-Stream on Zynq UltraScale+    & Vivado bitstream + PYNQ driver \\
2 & KV / TIU / MHC blocks land & Same AXI, more blocks visible              & Larger Vivado project \\
3 & post-tape-out (2027+)    & PCIe-attached custom chip                    & TSMC 16FFC silicon + Linux driver \\
\bottomrule
\end{tabular}
\end{table}

%==============================================================================
\section{Open Questions for the Compiler Team}
\label{sec:open}

These are in flight; the answers shape the v0.2 reference models.

\begin{enumerate}
\item \textbf{Compilation granularity.} Does your compiler emit
      individual matmuls plus a precision-gate op, or does it expect
      a fused \texttt{flash\_attention(Q, K, V)} op? Affects whether
      we add a fused-attention reference operation.
\item \textbf{FP16 bit-exactness.} Do you require bit-exact parity
      with our hardware's eventual FP16 rounding, or is
      numerically-close sufficient? See
      \texttt{mac\_array\_design.pdf} \S\ref{sec:fp16}.
\item \textbf{Async issue + token model.} Do you want non-blocking
      matmul issue today (we would add it to v0.2 now), or wait
      until the RTL exists?
\item \textbf{Cost-model precision.} Are placeholder cycle counts
      enough for your scheduler, or do you need post-synthesis
      numbers?
\item \textbf{\texttt{THRESHOLD} as a runtime register.} See ISA spec
      \S9, open question 1.
\end{enumerate}

We will track answers in the design docs and tag each decision as
\textbf{frozen} or \textbf{open} so the compiler team knows what
they can build against.

%==============================================================================
\section{Verification Summary}
\label{sec:verif}

\begin{table}[ht]
\centering
\small
\caption{Current verification status, both blocks.}
\label{tab:verif}
\begin{tabular}{@{}l l c@{}}
\toprule
Block                   & Test                                              & Status \\
\midrule
Precision Controller    & 143/143 RTL replay tiles, bit-exact (C++ + Py)    & pass \\
Precision Controller    & Stateful = batched = stateless                    & pass \\
Precision Controller    & extern "C" = C++ class                            & pass \\
Precision Controller    & Canonical edges (spike $=$ 10 / 11)               & pass \\
MAC Array               & INT8 matmul = naive int64 reference               & pass \\
MAC Array               & FP16 matmul = naive double reference (within tol) & pass \\
MAC Array               & Edge cases (zero, identity, signed)               & pass \\
MAC Array               & extern "C" = C++ class                            & pass \\
MAC Array               & FP16 round-trip helpers                           & pass \\
End-to-end              & Integration example routes match ground truth     & 82/18 vs 82/18 \\
\bottomrule
\end{tabular}
\end{table}

%==============================================================================
\section{Next Blocks in the Queue}
\label{sec:next}

In order of priority for the compiler co-design effort:

\begin{enumerate}
\item \textbf{Token Importance Unit.} Accumulator + classifier;
      $\sim$3--4 days once the keep/demote/evict threshold policy
      is decided.
\item \textbf{KV Cache Engine.} Codec decided --- \textbf{ChannelQuant}
      (per-channel INT4 keys / per-token INT4 values + FP16 outlier lane);
      RTL complete through Sky130 sign-off in the \texttt{kv-cache-engine} repo.
\item \textbf{Memory Hierarchy Controller.} Thin allocator first,
      refine when L1 / L2 / off-chip LPDDR5 behavior matters; $\sim$1 week.
\end{enumerate}

Each follows the same three-language template established by blocks
1 and 2.

\vspace{0.5in}
\noindent\rule{\textwidth}{0.4pt}\\
\noindent\small\textit{This document is the canonical entry point
for a compiler engineer working with LonghornSilicon. It is
versioned alongside the reference models in
\texttt{LonghornSilicon/attention-compute-unit} and updated
as each new block lands in software.}

\end{document}
